\chapter{Μεθοδολογία}
\label{ch:methods}

\section{Αλγόριθμοι Μηχανικής Μάθησης}
\label{sec:algorithms}

\subsection{LightGBM}

Το LightGBM~\cite{ke2017lgbm} αναπτύσσει τα δέντρα leaf-wise, με ταχύτερη σύγκλιση και
ελαφρώς υψηλότερη ευαισθησία σε υπερπροσαρμογή έναντι της level-wise ανάπτυξης
(§\ref{sec:ml_algorithms}, Σχ.~\ref{fig:leafwise_levelwise})· εδώ η υπερπροσαρμογή
ελέγχεται μέσω \texttt{num\_leaves}, \texttt{min\_child\_samples} και \texttt{subsample}.
Χρησιμοποιείται σε όλες τις στρατηγικές (TF, Rec, MIMO, Direct) και προτιμάται για τις
TF/Rec λόγω ταχύτητας εκπαίδευσης: ένα μοντέλο με $n \approx 1.000$ δέντρα εκπαιδεύεται σε
$\sim$1-2 λεπτά σε CPU, επιτρέποντας δεκάδες επαναλήψεις Optuna.


\subsection{XGBoost}

Το XGBoost~\cite{chen2016xgboost} χτίζει τα δέντρα level-wise, παράγοντας πιο συμμετρικές
δομές (§\ref{sec:ml_algorithms}). Εκπαιδεύεται πιο αργά από το LightGBM, αλλά εμπειρικά
παράγει πιο σταθερά αποτελέσματα σε δεδομένα υψηλής μεταβλητότητας - ιδιότητα που
αναδεικνύεται στην αξιολόγηση Q4 2025 και Q1 2026.


\subsection{Random Forest}

Το Random Forest~\cite{breiman2001rf} ακολουθεί parallel bagging - ανεξάρτητα δέντρα σε
τυχαία υποσύνολα δεδομένων και χαρακτηριστικών, με μέσο όρο των προβλέψεων - αντί για το
sequential boosting των LightGBM/XGB (§\ref{sec:ml_algorithms}). Η ανεξαρτησία των δέντρων
το καθιστά φυσικά ανθεκτικό σε ακραίες τιμές και έντονες διακυμάνσεις, όπως επαληθεύεται
στην περίοδο υψηλής μεταβλητότητας του Ιανουαρίου 2026.

Σε περιβάλλον Windows η παράλληλη εκτέλεση (multiprocessing) προκαλεί προβλήματα σταθερότητας· για τον λόγο αυτό το RF εκπαιδεύεται μονονηματικά (\texttt{n\_jobs=1}), ενώ η επιλογή \texttt{max\_features="sqrt"} μειώνει τον χρόνο εκπαίδευσης περίπου κατά $9\times$.


\subsection{Support Vector Regression}

Το SVR~\cite{smola2004} εκπαιδεύεται με πυρήνα RBF και $\varepsilon$-ανεκτική απώλεια,
ελαχιστοποιώντας τη νόρμα του μοντέλου (εξ.~\ref{eq:svr}-\ref{eq:svr_rbf},
§\ref{subsec:svr}). Λόγω της κυβικής επεκτασιμότητας ($O(n^3)$), εκπαιδεύεται σε υποσύνολο
5.000 τυχαίων δειγμάτων με τις εισόδους κανονικοποιημένες μέσω StandardScaler· ο
περιορισμός αυτός εξηγεί τη γενικά υποδεέστερη απόδοσή του σε μεγάλες δοκιμαστικές περιόδους.

\paragraph{Μεροληψία υπό μεταβολή καθεστώτος.} Πέρα από το υπολογιστικό κόστος, ο
πυρήνας ακτινικής βάσης (RBF kernel) $k(\mathbf{z}, \mathbf{z}') = \exp(-\gamma\,
\lVert \mathbf{z} - \mathbf{z}' \rVert^2)$ εισάγει μια δομική αδυναμία γενίκευσης εκτός
του εύρους της κατανομής εκπαίδευσης. Η πρόβλεψη γράφεται ως σταθμισμένο άθροισμα
τοπικών αποκρίσεων γύρω από τα διανύσματα υποστήριξης,
\begin{equation}
  f(\mathbf{z}) = \sum_{i \in \mathcal{S}} \alpha_i\, k(\mathbf{z}, \mathbf{z}_i) + b,
  \label{eq:svr_collapse}
\end{equation}
όπου $\mathcal{S}$ το σύνολο των διανυσμάτων υποστήριξης. Επειδή ο RBF είναι τοπικός,
για κάθε δοκιμαστικό σημείο $\mathbf{z}^\ast$ που απέχει σημαντικά από όλα τα
$\mathbf{z}_i$ ισχύει $k(\mathbf{z}^\ast, \mathbf{z}_i) \to 0$, οπότε $f(\mathbf{z}^\ast)
\to b$. Η πρόβλεψη καταρρέει προς τον σταθερό όρο $b$, ο οποίος είναι αγκυρωμένος κοντά
στη μέση τιμή των στόχων εκπαίδευσης. Όταν συμβεί έντονη μεταβολή καθεστώτος (regime
shift) - όπως η πτώση της μέσης τιμής από $108{,}7$ σε $78{,}3$~€/MWh τον
Φεβρουάριο 2026 - το νέο επίπεδο του στόχου βρίσκεται εκτός του φερόμενου εύρους των
διανυσμάτων υποστήριξης. Το μοντέλο αδυνατεί να παρεκβάλει και παράγει συστηματική
μεροληψία (bias), προβλέποντας σταθερά κοντά στο παρωχημένο μέσο επίπεδο. Η
$\varepsilon$-ανεκτική απώλεια επιτείνει το φαινόμενο, καθώς κλειδώνει τα διανύσματα
υποστήριξης στο κέντρο της αρχικής κατανομής και αδρανοποιεί τη λεπτή προσαρμογή στις
ακραίες περιοχές. Αυτό θεμελιώνει θεωρητικά τη μειωμένη ανθεκτικότητα του SVR στις
περιόδους μεταβολής καθεστώτος.


\subsection{Multi-Layer Perceptron}

Το MLP υλοποιείται σε PyTorch με αρχιτεκτονική (512, 256, 128) νευρώνων για τιμή και
(256, 128) για φορτίο, ενεργοποίηση ReLU, Dropout και BatchNormalization. Η εκπαίδευση
γίνεται με Adam optimizer (learning rate $5 \times 10^{-4}$) και early stopping (υπομονή
20 εποχών, παρακολούθηση validation loss). Για τη βελτιστοποίηση υπερπαραμέτρων MLP
χρησιμοποιείται ξεχωριστό Optuna study (τυνάρει hidden layers, dropout, learning rate,
weight decay, batch size).


\section{Στρατηγικές Πρόβλεψης Ορίζοντα 24 Ωρών}
\label{sec:strategies}

Οι τέσσερις στρατηγικές έχουν κοινό σκοπό: να προβλέψουν τις επόμενες 24 ώρες της
τιμής ή του φορτίου. Διαφέρουν στον τρόπο που χειρίζονται την εξάρτηση μεταξύ των ωρών
(βλ.\ Σχήμα~\ref{fig:strategies} στην §~\ref{sec:multistep}).
Ο Πίνακας~\ref{tab:strategies_summary} συνοψίζει τα κυριότερα χαρακτηριστικά τους.

Ο εβδομαδιαίος ορίζοντας $H=168$ αξιολογήθηκε επίσης κατά την έρευνα, αλλά δεν κρατήθηκε ως
κύριος άξονας παρουσίασης: πολλαπλασιάζει τον όγκο των αποτελεσμάτων χωρίς να μεταβάλλει τα
συμπεράσματα, ενώ το επιχειρησιακό σενάριο της ημερήσιας δημοπρασίας DAM αντιστοιχεί στο
$H=24$. Επιλεγμένα αποτελέσματα $H=168$ παρατίθενται μόνο όπου αναδεικνύουν τη συσσώρευση
σφάλματος της αναδρομικής (Πίν.~\ref{tab:conc_summary}).

\begin{table}[H]
\centering
\footnotesize
\caption{Συγκριτική παρουσίαση των τεσσάρων στρατηγικών πρόβλεψης $H=24$~ωρών.}
\label{tab:strategies_summary}
\begin{tabular}{lllll}
\toprule
Στρατηγική & Lagged είσοδοι & Error accum. & Μοντέλα & Χαρακτηρισμός \\
\midrule
TF        & Πραγματικές τιμές & Μηδέν       & 1           & Άνω φράγμα (oracle) \\
Rec       & Προβλέψεις        & Σταδιακό    & 1           & Αυτοπαλίνδρομη \\
MIMO      & Παρελθοντικές υστερήσεις & Μηδέν & 1 (multi-out) & Ταυτόχρονη πρόβλεψη \\
Direct    & Παρελθοντικές υστερήσεις & Μηδέν & $H=24$      & Ανεξάρτητα μοντέλα ανά ώρα \\
\bottomrule
\end{tabular}
\end{table}

\paragraph{Κοινός συμβολισμός.} Έστω $t$ η τελευταία διαθέσιμη ώρα (origin), που
αντιστοιχεί στο τέλος της ημέρας $D-1$, και $h \in \{1, \ldots, 24\}$ ο ορίζοντας
πρόβλεψης. Ορίζουμε το διάνυσμα εξωγενών χαρακτηριστικών $\mathbf{x}_\tau \in
\mathbb{R}^{m}$ (ημερολογιακά χαρακτηριστικά, τιμή αερίου, τιμή $\mathrm{CO_2}$,
ημερήσια πρόβλεψη φορτίου), όλα γνωστά ήδη στο κλείσιμο πύλης $D-1$, και το διάνυσμα
καθυστερημένων τιμών στόχου $\boldsymbol{\ell}_\tau = \left[y_{\tau-1}, y_{\tau-2},
\ldots, y_{\tau-168}\right]$. Το διάνυσμα εισόδου σχηματίζεται ως η συνένωση
$\mathbf{z}_\tau = \left[\boldsymbol{\ell}_\tau,\, \mathbf{x}_\tau\right]$. Η ουσιαστική
διαφορά των τεσσάρων στρατηγικών συνίσταται \emph{αποκλειστικά} στο ποια τιμή
τροφοδοτεί κάθε συνιστώσα του $\boldsymbol{\ell}$ όταν ο δείκτης υστέρησης πέφτει
\emph{εντός} του υπό πρόβλεψη 24ώρου, δηλαδή για $\tau - k > t$.


\subsection{Teacher-Forcing (TF)}

Το TF μοντέλο εκπαιδεύεται ως κλασικό μοντέλο παλινδρόμησης μίας εξόδου: για κάθε
παρατήρηση $t$, ο στόχος είναι $y_t$ και τα χαρακτηριστικά περιλαμβάνουν τις
καθυστερημένες τιμές $y_{t-1}, y_{t-2}, \ldots, y_{t-168}$. Κατά την αξιολόγηση, κάθε
συνιστώσα υστέρησης - συμπεριλαμβανομένων των ενδοημερήσιων ($t+h-k > t$) -
τροφοδοτείται με την \emph{πραγματική} παρατήρηση:

\begin{equation}
  \mathbf{z}^{\mathrm{TF}}_{t+h} = \big[\,\underbrace{y_{t+h-1}, y_{t+h-2}, \ldots,
  y_{t+h-168}}_{\text{πραγματικές υστερήσεις}},\; \mathbf{x}_{t+h}\,\big], \qquad
  \hat{y}_{t+h} = f\!\big(\mathbf{z}^{\mathrm{TF}}_{t+h}\big),
  \quad h = 1, \ldots, 24
  \label{eq:tf}
\end{equation}

Κάθε ώρα αξιολογείται μεμονωμένα, χωρίς συσσώρευση σφάλματος. Ωστόσο, η
αξιολόγηση Teacher-Forcing τροφοδοτεί σε κάθε βήμα την \emph{πραγματική}
καθυστερημένη τιμή του στόχου: για την πρόβλεψη της ώρας $t+h$ χρησιμοποιείται η
πραγματική $y_{t+h-1}$. Για ορίζοντες $h>1$ εντός της ημέρας, η τιμή αυτή
αντιστοιχεί σε ενδιάμεσες ώρες που \emph{δεν} έχουν ακόμη πραγματοποιηθεί τη
στιγμή υποβολής προσφορών (D-1, $\sim$12:00 CET). Συνεπώς, το TF λειτουργεί ως ιδεατό
μοντέλο αναφοράς (oracle) που ορίζει το άνω φράγμα της απόδοσης: τροφοδοτείται με τις
πραγματικές ιστορικές τιμές σε κάθε βήμα, ποσοτικοποιώντας τη μέγιστη εφικτή ακρίβεια. Χρησιμεύει
ως μέτρο του κόστους που επιφέρει η απουσία πρόσβασης στις πραγματικές ενδιάμεσες τιμές. Οι
στρατηγικές Recursive, MIMO και Direct προβλέπουν ολόκληρο το 24ωρο από την τελευταία
διαθέσιμη ώρα, χωρίς πρόσβαση σε ενδιάμεσες πραγματικές τιμές.


\subsection{Recursive (Rec)}

Η Rec στρατηγική χρησιμοποιεί το ίδιο μοντέλο μίας εξόδου, αλλά αντικαθιστά κάθε
ενδοημερήσια υστέρηση με την αντίστοιχη πρόβλεψη που έχει ήδη παραχθεί. Ορίζουμε τον
τελεστή διαθεσιμότητας
\begin{equation}
  \tilde{y}_\tau =
  \begin{cases}
    y_\tau, & \tau \le t \quad \text{(παρατηρηθείσα ιστορία)},\\[2pt]
    \hat{y}_\tau, & \tau > t \quad \text{(πρόβλεψη προηγούμενου βήματος)},
  \end{cases}
  \label{eq:rec_availability}
\end{equation}
οπότε το διάνυσμα εισόδου και η πρόβλεψη σχηματίζονται αναδρομικά ως
\begin{equation}
  \mathbf{z}^{\mathrm{Rec}}_{t+h} = \big[\,\tilde{y}_{t+h-1}, \tilde{y}_{t+h-2}, \ldots,
  \tilde{y}_{t+h-168},\; \mathbf{x}_{t+h}\,\big], \qquad
  \hat{y}_{t+h} = f\!\big(\mathbf{z}^{\mathrm{Rec}}_{t+h}\big),
  \quad h = 1, \ldots, 24
  \label{eq:recursive_methods}
\end{equation}
όπου σε κάθε βήμα η νέα πρόβλεψη $\hat{y}_{t+h}$ προστίθεται στο σύνολο των διαθέσιμων
τιμών και τροφοδοτεί τα επόμενα βήματα ($\tilde{y}_{t+h} \leftarrow \hat{y}_{t+h}$).

Η αναδρομή αυτή εισάγει σφάλμα σε κάθε βήμα. Για να αντιμετωπιστεί αυτό, εφαρμόζεται
Scheduled Sampling κατά την εκπαίδευση (βλ.\ Ενότητα~\ref{sec:scheduled_sampling}),
ώστε το μοντέλο να εκπαιδευτεί με αναφορά τόσο σε πραγματικές όσο και σε θορυβώδεις
τιμές.


\subsection{MIMO (Multi-Input Multi-Output)}

Το MIMO μοντέλο σχηματίζει ένα \emph{μοναδικό} διάνυσμα εισόδου από το origin $t$,
χρησιμοποιώντας αποκλειστικά υστερήσεις $\boldsymbol{\ell}_t$ με δείκτη $\le t$ (καμία
ενδοημερήσια συνιστώσα), και απεικονίζει το διάνυσμα αυτό απευθείας σε ολόκληρο τον
ορίζοντα 24 εξόδων:

\begin{equation}
  \mathbf{z}^{\mathrm{MIMO}}_{t} = \big[\,y_{t}, y_{t-1}, \ldots, y_{t-168},\;
  \mathbf{x}_{t}\,\big], \qquad
  \big[\hat{y}_{t+1}, \hat{y}_{t+2}, \ldots, \hat{y}_{t+24}\big]^{\!\top}
  = F\!\big(\mathbf{z}^{\mathrm{MIMO}}_{t}\big)
  \label{eq:mimo_methods}
\end{equation}

Καθώς το $\mathbf{z}^{\mathrm{MIMO}}_{t}$ παραμένει σταθερό για όλους τους ορίζοντες
$h$, δεν υπάρχει αναδρομική τροφοδότηση ($\tilde{y}_\tau$ με $\tau > t$ δεν εμφανίζεται
ποτέ) και επομένως κανένας μηχανισμός συσσώρευσης σφάλματος.

Στην υλοποίηση χρησιμοποιείται \texttt{MultiOutputRegressor} (για LightGBM, RF, SVR,
MLP) ή το XGBoost API με διαστάσεις εξόδου. Δεδομένου ότι δεν υπάρχει αναδρομική
τροφοδότηση, η ποιότητα των χαρακτηριστικών δεν υποβαθμίζεται για τις απομακρυσμένες
ώρες. Τα dense lags βοηθούν εδώ: τα $y_{\text{lag4}}, \ldots, y_{\text{lag23}}$
αναπαριστούν καθαρή ιστορία (όχι προβλέψεις) και παρέχουν πλούσια ενδοημερήσια
πληροφορία.

\paragraph{Two-stage MIMO για τιμή.} Για τα μοντέλα RF, SVR και MLP η πρόβλεψη τιμής
ακολουθεί δύο στάδια: (1) εκπαίδευση μοντέλου MIMO φορτίου $H=24$ (LGBM ή XGBoost),
(2) χρήση των προβλεπόμενων τιμών φορτίου ως επιπλέον 24 χαρακτηριστικά
($\hat{y}^{\text{load}}_{t+1}, \ldots, \hat{y}^{\text{load}}_{t+24}$) στο μοντέλο τιμής
(176 χαρακτηριστικά συνολικά). Η αλυσίδα αυτή εκμεταλλεύεται τη δομική σχέση
τιμής-φορτίου χωρίς να εισάγει leakage.


\subsection{Direct Multistep}

Η Direct στρατηγική εκπαιδεύει $H = 24$ ανεξάρτητα μοντέλα, ένα ανά ορίζοντα. Όπως και
στο MIMO, το διάνυσμα υστερήσεων αντλείται αποκλειστικά από το origin $t$ (δείκτες
$\le t$), ενώ τα εξωγενή χαρακτηριστικά λαμβάνονται για τη συγκεκριμένη ώρα-στόχο
$t+h$:

\begin{equation}
  \mathbf{z}^{\mathrm{Dir}}_{t,h} = \big[\,y_{t}, y_{t-1}, \ldots, y_{t-168},\;
  \mathbf{x}_{t+h}\,\big], \qquad
  \hat{y}_{t+h} = f_h\!\big(\mathbf{z}^{\mathrm{Dir}}_{t,h}\big),
  \quad h = 1, \ldots, 24
  \label{eq:direct_methods}
\end{equation}

Κάθε μοντέλο $f_h$ εκπαιδεύεται αποκλειστικά στον υπολογισμό της τιμής $h$ ώρες αργότερα.
Απαιτεί $H\times$ παραπάνω χρόνο εκπαίδευσης, αλλά εξαλείφει τη συσσώρευση
σφάλματος. Στην υλοποίηση αντιστοιχεί στο LightGBM Direct (24 ανεξάρτητα μοντέλα).


\section{Scheduled Sampling}
\label{sec:scheduled_sampling_methods}

Το Scheduled Sampling~\cite{bengio2015scheduled} προσαρμόζεται ως εξής για αλγόριθμους
gradient boosting:

\begin{enumerate}
  \item Αρχικοποίηση του μοντέλου με $\varepsilon_0 = 0{,}10$ (10\% των lagged
  χαρακτηριστικών αντικαθίστανται από pseudo-προβλέψεις).

  \item Εκτέλεση ημερήσιου 24ωρου recursive loop για κάθε ημέρα εκπαίδευσης: για
  κάθε ώρα $h \in \{1,\ldots,24\}$, αποφασίζεται τυχαία (με πιθανότητα $\varepsilon$)
  αν το $y_{\text{lag}_k}$ θα είναι η πραγματική τιμή ή η προβλεπόμενη.

  \item Επανάληψη για $N_{\text{iter}} = 3$ επαναλήψεις, αυξάνοντας το $\varepsilon$
  γραμμικά:
\end{enumerate}

\begin{equation}
  \varepsilon_t = \varepsilon_{\min} + \frac{t}{N_{\text{iter}}-1}
  \left(\varepsilon_{\max} - \varepsilon_{\min}\right), \quad
  \varepsilon_{\min}=0{,}10,\; \varepsilon_{\max}=0{,}40
  \label{eq:ss_schedule}
\end{equation}

\begin{figure}[htbp]
  \centering
  \includegraphics[width=0.88\textwidth]{fig_scheduled_sampling}
  \caption{Πρόγραμμα Scheduled Sampling: η πιθανότητα $\varepsilon$ χρήσης
    προβλεπόμενης τιμής αντί πραγματικής αυξάνει γραμμικά από $0{,}10$
    έως $0{,}40$ σε τρία περάσματα εκπαίδευσης, υλοποιώντας τη μετάβαση
    teacher-forcing $\rightarrow$ autoregressive (βλ.\ Εξ.~\ref{eq:ss_schedule}).}
  \label{fig:ss_schedule}
\end{figure}

Αυτή η διαδικασία κάνει το μοντέλο ανθεκτικό στον θόρυβο των lagged εισόδων κατά τη
recursive χρήση. Ειδικά για το φορτίο, το Scheduled Sampling βελτιώνει σημαντικά τα
αποτελέσματα (βλ.\ Κεφάλαιο~\ref{ch:experiments}).


\section{Βελτιστοποίηση Υπερπαραμέτρων με Optuna}
\label{sec:optuna}

Η βελτιστοποίηση γίνεται ανεξάρτητα για κάθε αλγόριθμο και κάθε εργασία, χρησιμοποιώντας
το Optuna με τον αλγόριθμο TPE. Ο Πίνακας~\ref{tab:optuna_params} συνοψίζει τα εύρη
αναζήτησης και τις βέλτιστες τιμές.

\begin{table}[H]
\centering
\small
\caption{Εύρη αναζήτησης υπερπαραμέτρων (Optuna) και βέλτιστες τιμές για το καλύτερο
μοντέλο κάθε αλγόριθμου (εργασία τιμής, TF/Rec στρατηγική).}
\label{tab:optuna_params}
\begin{tabular}{llll}
\toprule
Αλγόριθμος & Υπερπαράμετρος & Εύρος & Βέλτιστη τιμή \\
\midrule
\multirow{5}{*}{LightGBM}
  & n\_estimators      & [200, 2000]               & 900   \\
  & learning\_rate     & [0.005, 0.15]             & 0.027 \\
  & num\_leaves        & [15, 255]                 & 26    \\
  & min\_child\_samples & [10, 100]                & 54    \\
  & subsample          & [0.5, 1.0]                & 0.925 \\
\midrule
\multirow{4}{*}{XGBoost}
  & n\_estimators      & [200, 1500]               & 405   \\
  & learning\_rate     & [0.01, 0.30]              & 0.030 \\
  & max\_depth         & [3, 12]                   & 12    \\
  & subsample          & [0.5, 1.0]                & 0.85  \\
\midrule
\multirow{3}{*}{MLP}
  & hidden\_layers     & \{(128,64),(256,128),(512,256,128)\} & (512,256,128) \\
  & dropout            & [0.0, 0.5]                & 0.37  \\
  & learning\_rate     & [$10^{-4}$, $10^{-2}$]   & $9{,}6 \times 10^{-4}$ \\
\midrule
\multirow{2}{*}{SVR}
  & C                  & [$10^{-1}$, $10^3$] (log) & 50.0  \\
  & epsilon            & [$10^{-3}$, 1.0] (log)    & 0.05  \\
\midrule
\multirow{2}{*}{LGBM (MIMO)}
  & n\_estimators      & [300, 2000]               & 900   \\
  & num\_leaves        & [15, 127]                 & 26    \\
\bottomrule
\end{tabular}
\end{table}


\paragraph{TF/Rec Βελτιστοποίηση.}
Για τα TF και Rec μοντέλα χρησιμοποιείται \textbf{ημερήσιο rolling-origin CV}: κάθε fold
αξιολογεί το μοντέλο σε ένα 24ωρο παράθυρο, ακολουθώντας αυστηρά χρονολογική σειρά
ώστε να μην υπάρχει διαρροή από το μέλλον (no look-ahead). Το παράθυρο εκπαίδευσης
επεκτείνεται σωρευτικά: κάθε νέο fold προσθέτει τα δεδομένα του προηγούμενου 24ώρου
στο training set. Αυτό προσομοιώνει πιστά τον πραγματικό τρόπο ανάπτυξης του μοντέλου
στην αγορά. Για τα Rec μοντέλα η αξιολόγηση στο CV εκτελείται επίσης αναδρομικά
(recursive prediction, όχι teacher-forcing), ώστε η CV-απόδοση να αντικατοπτρίζει
πραγματικές συνθήκες χρήσης και να μην υπερεκτιμά λόγω διαρροής teacher-forcing.
Ο αριθμός των splits ορίζεται σε 3-5, καλύπτοντας τις τελευταίες εβδομάδες του
εκπαιδευτικού συνόλου.

\paragraph{MIMO/Direct Βελτιστοποίηση.}
Για τα MIMO και Direct μοντέλα η δομή του CV διαφέρει: αντί για 1-step-ahead αξιολόγηση,
κάθε fold αξιολογεί το μοντέλο σε πλήρη 24ωρα προβλεπτικά παράθυρα. Η συνάρτηση κόστους
του Optuna είναι το \textit{μέσο MAE} σε ολόκληρο τον ορίζοντα $H = 24$, ισοσταθμίζοντας
όλες τις ώρες. Αυτή η επιλογή είναι κρίσιμη: ένα μοντέλο που ελαχιστοποιεί μόνο το
σφάλμα για $h=1$ θα αποδίδει κακά για $h=20$-$24$. Εναλλακτικά, θα μπορούσε να
χρησιμοποιηθεί weighted MAE που να δίνει μεγαλύτερη βαρύτητα στους απόμακρους ορίζοντες,
αλλά η ισοσταθμία αποδείχθηκε ισοδύναμη στην πράξη.

\paragraph{Walk-forward Βελτιστοποίηση.}
Στο πλαίσιο της walk-forward μηνιαίας επανεκπαίδευσης (βλ.\ §~\ref{sec:walk_forward})
τίθεται το ερώτημα πόσο συχνά πρέπει να επαναλαμβάνεται η πλήρης βελτιστοποίηση. Η
ανεξάρτητη εκτέλεση Optuna (40-60 trials) κάθε μήνα θα ήταν πληρέστερη αλλά απαιτητική
υπολογιστικά. Η λύση που υιοθετείται είναι \textbf{CachedOptuna}: για τον Δεκέμβριο 2025
χρησιμοποιούνται προ-υπολογισμένες βέλτιστες υπερπαράμετροι από Optuna σε πλήρες
εκπαιδευτικό σύνολο (40-60 trials, cutoff Νοε.~2025). Για Ιανουάριο και Φεβρουάριο
2026 εκτελείται γρήγορο Optuna (15 trials, 2 splits) με ενημερωμένα δεδομένα, χωρίς
να ξεκινά από μηδέν: ο αλγόριθμος TPE αρχικοποιείται με τις cached παραμέτρους ως
σημείο εκκίνησης. Η προσέγγιση αυτή μειώνει τον χρόνο βελτιστοποίησης κατά $\sim$70\%
διατηρώντας παρεμφερή ποιότητα υπερπαραμέτρων.


\section{Ensembles}
\label{sec:ensembles}

Για κάθε στρατηγική κατασκευάζονται δύο σύνολα ensembles:

\textbf{Ensemble-Best3 (1/MAE weighted).} Επιλέγονται τα 3 καλύτερα μοντέλα βάσει
ιστορικού MAE και σταθμίζονται αντιστρόφως ανάλογα με το σφάλμα τους:

\begin{equation}
  \hat{y}^{\text{ens}}_t = \frac{\sum_{i=1}^{3} w_i \hat{y}^{(i)}_t}{\sum_{i=1}^{3} w_i},
  \quad w_i = \frac{1}{\mathrm{MAE}_i}
  \label{eq:ensemble}
\end{equation}

\textbf{Ensemble-All (1/MAE weighted).} Ίδια στάθμιση αλλά με όλα τα διαθέσιμα μοντέλα
(εξαιρείται μόνο το SVR για τιμή λόγω κακής σχετικής απόδοσης).

\textbf{Αποφυγή διαρροής στη σύνθεση.} Η επιλογή των τριών μοντέλων και τα βάρη
$w_i$ προκύπτουν αποκλειστικά από \emph{validation} σφάλμα (στο walk-forward,
το MAE του \emph{προηγούμενου} μήνα) και ποτέ από το σύνολο δοκιμής·
διαφορετικά η σύνθεση θα ενσωμάτωνε look-ahead πληροφορία (επιλογή και στάθμιση
των μοντέλων με γνώση της μελλοντικής τους επίδοσης). Υπό αυτό το αυστηρό σχήμα, το
Ensemble-Best3 \emph{δεν} υπερτερεί συστηματικά του καλύτερου μεμονωμένου μοντέλου:
για την τιμή το leak-free Ensemble δίνει $\approx 11{,}6$~€/MWh έναντι
$10{,}46$~€/MWh του μεμονωμένου TF-LGBM-MR, ενώ για το φορτίο τα δύο είναι
ισοδύναμα ($\approx 69$~MW). Για τον λόγο αυτό, ως κύριο μοντέλο αναφοράς
χρησιμοποιείται το μεμονωμένο TF-LGBM-MR και το ensemble παρουσιάζεται ως
ισοδύναμη, ανθεκτική εναλλακτική.


\section{Μετρικές Αξιολόγησης}
\label{sec:metrics}

Η αξιολόγηση κάθε μοντέλου βασίζεται σε τρεις μετρικές:

\begin{equation}
  \mathrm{MAE} = \frac{1}{T}\sum_{t=1}^{T} |y_t - \hat{y}_t|
  \label{eq:mae}
\end{equation}

Κύρια μετρική αξιολόγησης. Είναι ανθεκτική σε ακραίες τιμές και εύκολα ερμηνεύσιμη
(€/MWh ή MW).

\begin{equation}
  \mathrm{RMSE} = \sqrt{\frac{1}{T}\sum_{t=1}^{T} (y_t - \hat{y}_t)^2}
  \label{eq:rmse}
\end{equation}

Τιμωρεί δυσανάλογα τα μεγάλα σφάλματα, επομένως είναι ευαίσθητη στα price spikes.

\begin{equation}
  \mathrm{sMAPE} = \frac{100}{T}\sum_{t=1}^{T}\frac{|y_t - \hat{y}_t|}{|y_t| + |\hat{y}_t|}
  \label{eq:smape}
\end{equation}

Κανονικοποιημένη εκδοχή του ποσοστιαίου σφάλματος, συμμετρική ως προς υπερ- και
υπο-εκτίμηση. Χρήσιμη για σύγκριση μεταξύ περιόδων διαφορετικού επιπέδου τιμής. Η
κύρια μετρική για την κατάταξη είναι το MAE.


\section{Βασικές Αναφορές (Baselines)}
\label{sec:baselines}

Στην προκειμένη εργασία χρησιμοποιούνται τέσσερις αναφορές (baselines):

\begin{itemize}
  \item \textbf{Naive-1:} $\hat{y}_{t+1} = y_t$ (επανάληψη τελευταίας τιμής),
  \item \textbf{Naive-24:} $\hat{y}_{t+h} = y_{t+h-24}$ (ίδια ώρα χθεσινής ημέρας),
  \item \textbf{Naive-168:} $\hat{y}_{t+h} = y_{t+h-168}$ (ίδια ώρα προηγούμενης εβδομάδας),
  \item \textbf{Seasonal Profile:} $\hat{y}_{t+h} = \operatorname{mean}\{\, y_\tau : \tau \in \mathcal{T}_{\mathrm{tr}},\ (\mathrm{hod}_\tau, \mathrm{dow}_\tau) = (\mathrm{hod}_{t+h}, \mathrm{dow}_{t+h}) \,\}$, ο μέσος όρος του εκπαιδευτικού συνόλου $\mathcal{T}_{\mathrm{tr}}$ ανά (ώρα ημέρας, ημέρα εβδομάδας).
\end{itemize}

Ένα μοντέλο θεωρείται αξιόλογο μόνο αν υπερτερεί τουλάχιστον του Naive-1. Για τη
στρατηγική MIMO, ο πιο φυσικός baseline είναι το Naive-24, καθώς αμφότερα προβλέπουν
ολόκληρη 24ωρη ημέρα από τα ίδια ιστορικά δεδομένα.
